\documentclass[a4paper, 12pt]{article}
\usepackage[utf8]{inputenc}
\usepackage[margin=1in]{geometry}
\usepackage{setspace}
\usepackage{xcolor}
\usepackage{hyperref}
\usepackage{lipsum}
\usepackage{fancyhdr}
\usepackage[para]{footmisc}

\renewcommand{\headrulewidth}{0pt}
\fancyfoot[R]{\thepage}
\singlespacing
\usepackage{titlesec}
\titleformat{\section}
  {\normalfont\large\bfseries}{\thesection}{1em}{}
\titleformat{\subsection}
  {\normalfont\normalsize\bfseries}{\thesubsection}{1em}{}
\titleformat{\subsubsection}
  {\normalfont\small\bfseries}{\thesubsubsection}{1em}{}
\titlespacing*{\section}{0pt}{1.5ex plus 1ex minus .2ex}{1ex plus .2ex}
\titlespacing*{\subsection}{0pt}{1.5ex plus 1ex minus .2ex}{0.8ex plus .2ex}
\titlespacing*{\subsubsection}{0pt}{1ex plus 1ex minus .2ex}{0.5ex plus .2ex}
\renewcommand{\thesubsection}{\arabic{subsection}}

\pagestyle{fancy}
\fancyhf{}
\fancyhead[L]{IK2221 VT26}
\fancyhead[R]{Lorenzo Deflorian (ldef@kth.se)}
%\fancyfoot[R]{\thepage}

\begin{document}

\section*{The Click Modular Router}
\subsection{Paper summary}
% Paper summary (for E level)
%% Problem being tackled and why it's hard
%% Key idea behind the approach
%% A bit more detail about the approach
%% The way the idea is evaluated, plus the key findings

This paper argues that traditional routers are too rigid and not flexible enough for all the functions they are expected to handle today, like filtering, queueing, firewalls, tunneling, etc. Basically, everything is too ``closed'' and hard to modify.

To solve this, the authors introduce Click, which is a modular router architecture. Instead of having one big monolithic system, the router is built from small components called elements. Each element does a simple task (for example classification or queueing), and they are connected together in a graph that defines how packets move through the router.

An interesting part is that Click supports both push and pull models, so packets can either be pushed forward or requested by the next element. This makes the design more flexible, even if it's not always super intuitive at first. There is also something called flow-based context, which allows elements to find related ones automatically, and a configuration language to describe the router setup.

The authors show that this approach actually works by building a full IP router and then extending it with more advanced features like RED, fair queueing, and tunneling. So it's not just a theoretical idea.

In terms of performance, Click runs pretty well on PC hardware. They show that it can process around 333k packets per second, which is significantly faster than standard Linux routing on the same machine, even if the comparison is a bit specific. Overall, the main point is that you can have both flexibility and good performance, and not necessarily sacrifice one for the other.

\subsection{Contributions}
% Top 3 contributions of the paper  (extra, for C level)
%% Tie the contributions with impact.
%% The first paper to propose… | Makes it possible to… 
%% Presents a new technique for… | Dramatically improves performance of … 
%% Offers largely improved reliability… | Saves a large fraction of power for…

The first contribution is the idea of breaking down a router into small modular elements connected in a graph. This makes it much easier to extend or modify the router without redesigning everything from scratch, which is a big improvement compared to traditional designs.

The second contribution is the set of design features that make this model actually usable. In particular, push/pull connections, explicit queues, flow-based context, and the configuration language. These things give a lot of control over how packets are handled, even if they also make the system a bit harder to fully understand at the beginning.

The third contribution is the implementation and evaluation. The authors don't just propose the idea, but they build a real system and show that it performs well. The fact that it can be faster than standard Linux routing (in their setup) is important, because otherwise the modularity argument would be less convincing.

\subsection{Problems}
% Most glaring problem(s) with the work (up to 3) (extra, for A level)
%% Does the approach scale as the number of users increases?
%% Scaling: non-linear increase in resources | How does it work under a failure?
%% Does it miss some important cases? | Does it have too strong assumptions?
%% Does it have unrealistic deployment requirements?

One limitation is that Click works very well for packet processing, but not as well for more complex control logic. For example, protocols like BGP don't really fit nicely into this model, since they don't follow the same kind of packet flow. The paper mentions this, but it still feels like a gap in the design.

Another issue is configuration complexity. Even if Click makes routers more flexible, it also puts more responsibility on the user. You need to understand how to connect elements correctly, how push and pull interact, and how queues behave. If you get something wrong, it might still work but not in the way you expect, which is a bit risky.

Finally, the performance results are good, but they depend a lot on the specific setup (old Linux kernel, single-threaded system, etc). So while the results are convincing for that context, it's not completely clear how well this would scale to modern systems with higher speeds and more cores. This part is not really explored in the paper.

\end{document}